\section{Kuantor Eksistensial: Saksi Kebenaran dan Konjungsi}

Berbeda dari kuantor universal, kuantor eksistensial berfokus pada keberadaan minimal satu elemen yang memenuhi properti tertentu \cite{rosen2019}.

\subsection{Definisi Formal Kuantor Eksistensial}

\begin{definisi}[Kuantor Eksistensial]
Kuantifikasi eksistensial merubah fungsi proposisional $P(x)$ menjadi proposisi ``terdapat setidaknya satu (sekurang-kurangnya satu) elemen $x$ di dalam semesta pembicaraan $\mathcal{U}$ sedemikian sehingga $P(x)$ bernilai benar''. Notasinya ditulis:
\[ \exists x \, P(x) \]
Simbol $\exists$ asalnya dari huruf ``E'' terbalik (\textit{Exists}).
\end{definisi}

Dibaca:
\begin{itemize}
    \item ``Terdapat sebuah $x$ sedemikian rupa sehingga $P(x)$''
    \item ``Sekurang-kurangnya satu $x$ memenuhi $P(x)$''
    \item ``Ada $x$ dengan sifat $P(x)$''
\end{itemize}

Jika semesta pembicaraan $\mathcal{U}$ berbentuk himpunan terhingga $\mathcal{U} = \{a_1, a_2, \ldots, a_n\}$, kuantor eksistensial ekuivalen dengan disjungsi seluruh elemen:
\[ \exists x \, P(x) \equiv P(a_1) \lor P(a_2) \lor \ldots \lor P(a_n) \]

Karena ekuivalen dengan disjungsi, kalimat eksistensial bernilai \textbf{Salah} hanya jika \emph{semua} elemen domain gagal memenuhi $P(x)$. Sebaliknya, cukup ditemukan satu nilai $c$ dengan $P(c)$ bernilai benar (disebut \textit{witness} atau saksi eksistensi), maka keseluruhan proposisi bernilai benar.

\subsection{Kombinasi dengan Konjungsi ($\land$)}

Dalam translasi dari bahasa alami, frasa berbau eksistensial lazim diterjemahkan menjadi operator konjungsi logika.

\begin{contoh}
Translasi proposisi: ``Sebagian buah terasa asam.''
\begin{itemize}
    \item Jika domain $\mathcal{U}$ direstriksi pada himpunan buah-buahan saja, biarkan $A(x) \equiv \text{``}x \text{ terasa asam''}$. Notasinya ringkas:
    \[ \exists x \, A(x) \]
    
    \item Apabila semesta $\mathcal{U}$ dibiarkan amat luas mencakup segala benda di dapur, biarkan $B(x) \equiv \text{``}x \text{ adalah buah''}$. Proposisi menyatakan bahwa \textit{terdapat benda yang merupakan buah DAN secara bersamaan ia juga terasa asam}. Maka translasinya:
    \[ \exists x \, (B(x) \land A(x)) \]
\end{itemize}
\end{contoh}

\begin{catatan}
\textbf{Perhatian!} Kuantor eksistensial biasanya tidak dipasangkan dengan implikasi ($\rightarrow$) untuk pola ``sebagian A adalah B''. Bentuk $\exists x \, (B(x) \rightarrow A(x))$ mudah menjadi benar secara \textit{vacuous truth}: elemen yang bukan buah sudah cukup membuat implikasi benar. Untuk pola ``ada A yang B'', bentuk baku adalah $\exists x \, (B(x) \land A(x))$ \cite{wiki_quantifier_logic}.
\end{catatan}

\subsection{Contoh Matematis Eksistensial}

\begin{contoh}
Diberikan domain $\mathcal{U} = \text{Bilangan Real } (\mathbb{R})$.
\begin{enumerate}
    \item $\exists x \ (x^2 = 9)$ \\ \textbf{Nilai: Benar}. Salah satu \textit{witness} adalah $x = 3$, karena $3^2 = 9$.
    \item $\exists x \ (x = x + 1)$ \\ \textbf{Nilai: Salah}. Tidak ada bilangan riil yang sama dengan dirinya setelah ditambah 1.
    \item $\exists x \ (x^2 < 0)$ \\ \textbf{Nilai: Salah} pada domain $\mathbb{R}$. Jika domain diubah ke $\mathbb{C}$, terdapat nilai seperti $x=i$ dengan $i^2=-1$, sehingga klaim menjadi benar.
\end{enumerate}
\end{contoh}
